% !TeX encoding = UTF-8
% !TeX program = pdflatex
\documentclass[lettersize,journal]{IEEEtran}

% Mathematics
\usepackage{amsmath,amssymb,amsfonts}
\usepackage{mathtools}

% Tables
\usepackage{array}
\usepackage{tabularx}
\usepackage{booktabs}
\usepackage{multirow}
\usepackage{colortbl}

% Figures and graphics
\usepackage{graphicx}
\usepackage[caption=false,font=normalsize,labelfont=sf,textfont=sf]{subfig}
\usepackage{rotating}
\usepackage{stfloats}

% Algorithms
\usepackage{algorithm}
\usepackage{algorithmic}

% Symbols and utilities
\usepackage{textcomp}
\usepackage{pifont}
\usepackage{afterpage}
\usepackage{balance}
\usepackage{verbatim}
\usepackage{xcolor}
\usepackage{url}

% Citations and references
\usepackage{cite}

% Hyperlinks
\usepackage{hyperref}

\newcolumntype{Y}{>{\centering\arraybackslash}X}
\hypersetup{
    linktocpage=true,
    pdfborderstyle={/S/S/W 1},
    hyperindex=true,
    bookmarks=false,
    bookmarksopen=false,
    bookmarksnumbered=false
}

% Paths for manuscript assets.
\graphicspath{{Fig/}{Plot/}{Appendix/Fig/}}

% Hyphenation
\hyphenation{op-tical net-works semi-conduc-tor IEEE-Xplore}

% Compatibility fixes
\let\Bbbk\relax

\begin{document}

\title{Paper Title}

\author{Author~Name,
        Second~Author,
        and~Third~Author%
\thanks{Author Name is with Affiliation, City, Country (e-mail: author@example.com).}
\thanks{Manuscript received Month Day, Year; revised Month Day, Year.}}

\markboth{IEEE Transactions on Journal Name,~Vol.~XX, No.~XX, Month~Year}%
{Author \MakeLowercase{\textit{et al.}}: Paper Title}

\maketitle

\begin{abstract}
Write a concise summary of the problem, method, main results, and significance.
\end{abstract}

\begin{IEEEkeywords}
Keyword one, keyword two, keyword three, keyword four.
\end{IEEEkeywords}

\section{Introduction}
\IEEEPARstart{A}{I}-generated images have become increasingly diverse as generative models evolve from adversarial synthesis to diffusion and text-conditioned generation~\cite{goodfellow2014gan,karras2019stylegan,ho2020ddpm,dhariwal2021diffusion,rombach2022ldm}. Modern systems can synthesize realistic faces, objects, scenes, and artwork under a wide range of prompts and post-processing pipelines. This progress expands creative applications, but it also makes image authenticity harder to assess from pixels alone. AI-generated image detection therefore has to operate in a setting where both generators and visual content distributions change over time.

This paper studies image-level detection from RGB pixels. Given an image, the detector must assign a fake score without relying on generator internals, prompts, provenance metadata, or watermarks at test time. This setting is deliberately narrower than full media provenance, but it remains important because many shared images arrive without trusted side information. The practical question is not only whether a detector can recognize a known generator, but whether its score remains meaningful when the synthesis process and post-processing pipeline differ from those seen during development.

The central difficulty is cross-generator generalization. A detector trained on one source generator can perform well when target images contain similar artifacts, yet fail when the generator architecture, image content, resolution, compression, or degradation pipeline changes. Early cross-generator studies showed that ProGAN-trained classifiers can transfer to several unseen CNN-based generators when training and preprocessing are carefully designed~\cite{wang2020cnnspot}. Later work broadened the setting to vision-language feature spaces, diffusion models, commercial generators, challenging user-curated images, and real-world degradations~\cite{ojha2023univfd,zhu2023genimage,yan2025aide,li2025rrdataset}. These studies make generalization, rather than in-domain accuracy alone, the central requirement for deployable detection.

Existing detectors approach this requirement through different evidence sources. Spatial classifiers and texture methods learn visual artifacts left by synthesis pipelines~\cite{wang2020cnnspot,liu2020gramnet,zhong2023patchcraft}. Frequency and statistical methods search for traces in spectral, noise, color, bit-plane, or photographic regularities~\cite{frank2020frequency,liu2022realimages,karageorgiou2025spai,zhong2025beyond,wang2025lota}. Residual, gradient, and reconstruction methods expose generated images through model gradients, upsampling residuals, diffusion reconstruction, autoencoder reconstruction, or inversion behavior~\cite{tan2023lgrad,tan2024npr,wang2023dire,ricker2024aeroblade,luo2024lare2,cazenavette2024fakeinversion}. Semantic and foundation-model detectors use pretrained representations, such as CLIP, to improve transfer beyond the source generator~\cite{radford2021clip,ojha2023univfd,cozzolino2024clip}.

Recent top-conference work shows that the field is moving beyond a single artifact family. Bias-free training and information-bottleneck objectives reduce source-data shortcuts~\cite{guillaro2025biasfree,zhang2025vibnet}. Any-resolution spectral learning, bit-plane modeling, real-image distribution modeling, camera-metadata supervision, and color-correlation training seek low-level cues that are less dependent on one generator family~\cite{karageorgiou2025spai,wang2025lota,liu2022realimages,zhong2025beyond,zhong2026metadata,zhong2026color}. At the same time, hybrid detectors combine semantic, pixel-level, generative, or multi-expert cues to improve transfer~\cite{yan2025aide,cheng2025cospy,tan2026mcan,li2026synerdetect}. This progress suggests that generated-image detection is best viewed as a multi-prior problem, rather than as a search for one universal trace.

However, multi-prior detection introduces a second challenge. Low-level artifacts can be generator-specific or altered by resizing and compression; semantic features can transfer across content categories but may miss subtle synthesis traces; reconstruction scores can be informative but inherit assumptions from the inverse process, autoencoder, or diffusion prior that defines the residual. Combining these cues can improve robustness, but it also creates choices about score normalization, fusion weights, checkpoints, and thresholds. If those choices are tuned with target labels, a cross-generator benchmark becomes a validation set rather than a report-only test.

Evaluation protocol is therefore part of the technical problem. Benchmarks and sanity-check studies show that detector behavior can change across datasets, generator families, image degradations, challenge sets, and calibration choices~\cite{yan2023deepfakebench,zhu2023genimage,yan2025aide,li2025rrdataset,yang2026calibrated}. In this setting, target labels are often available during analysis and can unintentionally influence thresholds, fusion weights, score variants, or the subset of generators emphasized in a paper. For cross-generator detection, how a detector is selected can be as important as what detector is selected.

We propose FreqPRISM, a source-only multi-prior detector for AI-generated image detection. FreqPRISM integrates four complementary scores: a whole-image artifact score from explicit low-level statistics, a native-tile artifact score that preserves local high-resolution evidence, a semantic score from a frozen CLIP image encoder with a source-fitted linear probe, and a residual score from nearest-neighbor down-up reconstruction traces. The components are fused in logit space with constants selected from source data only, and the operating threshold is fixed at $0.50$ before target reporting. Target labels are used only after the detector is locked.

The resulting pipeline is designed to make prior integration auditable. The artifact prior anchors the detector in explicit low-level evidence; the tile pathway preserves localized high-resolution traces; the semantic prior supplies representation-level transfer; and the residual prior contributes a controlled correction signal. We evaluate the locked detector on AIGCDetectBenchmark, the 17-generator benchmark introduced by PatchCraft~\cite{zhong2023patchcraft}, analyze the role of each prior through component and artifact-family ablations, examine fusion sensitivity, and report benchmark transfer on UniversalFakeDetect and GenImage. These experiments test not only ranking quality, but also fixed-threshold behavior and calibration under domain shift.

Our contributions are summarized as follows:
\begin{itemize}
    \item We formulate FreqPRISM as a source-only prior-integration detector for cross-generator AI-generated image detection, with target labels reserved for report-only evaluation.
    \item We combine whole-image artifact, native-tile artifact, frozen-CLIP semantic, and nearest-neighbor down-up residual evidence into one fixed-threshold fake score.
    \item We define a source-only stress-calibrated logit fusion protocol that fixes fusion constants and the decision threshold before target evaluation.
    \item We provide component-level analyses on AIGCDetectBenchmark and external UniversalFakeDetect and GenImage reports, while explicitly identifying calibration and domain-shift boundaries.
\end{itemize}

\section{Related Work}
\label{sec:related_work}

AI-generated image detection has been approached through several evidence families, each making a different assumption about what should remain stable across generators. Spatial detectors rely on visual artifacts left by synthesis pipelines~\cite{wang2020cnnspot,liu2020gramnet,zhong2023patchcraft}; frequency and photographic-statistics methods search for signal-level regularities~\cite{frank2020frequency,liu2022realimages,zhong2025beyond}; residual and reconstruction methods use inverse or perturbation behavior as evidence~\cite{tan2023lgrad,tan2024npr,wang2023dire}; and semantic or hybrid detectors exploit transferable representation spaces~\cite{radford2021clip,ojha2023univfd,yan2025aide,cheng2025cospy}. We therefore group prior work into low-level forensic cues, residual and reconstruction cues, semantic and hybrid detectors, and evaluation protocols, emphasizing the evidence assumptions behind each line rather than benchmark ranking alone.

\subsection{Low-Level Artifact and Forensic Cues}
Early cross-generator studies showed that synthesis pipelines leave learnable spatial traces. CNNSpot demonstrated that a ProGAN-trained classifier can transfer to several unseen CNN-based generators when augmentation and preprocessing are carefully controlled~\cite{wang2020cnnspot}. This protocol remains influential, but it also exposes a weakness of purely supervised artifact learning: the detector can inherit source-generator and source-data biases rather than a stable notion of photographic authenticity.

Later spatial and texture-oriented detectors made the artifact hypothesis more structured. GramNet models global texture statistics for fake-face detection~\cite{liu2020gramnet}, PatchCraft emphasizes local texture patches and rich/poor texture contrast~\cite{zhong2023patchcraft}, and Fusing combines image-level and local evidence for generalized synthesized-image detection~\cite{ju2022fusing}. More recent architectures, including FatFormer and FerretNet, improve transfer by adapting transformer features or exploiting local-pixel dependency patterns~\cite{liu2024fatformer,liang2025ferretnet}. Other work targets the source-bias problem more directly through bias-free generation protocols or information-bottleneck objectives~\cite{guillaro2025biasfree,zhang2025vibnet}. These studies motivate an artifact prior, but they also argue against treating a single artifact classifier as sufficient for open-generator detection.

Frequency and photographic-statistics methods seek traces that are less semantic and potentially less tied to image content. FreDect showed that GAN images can contain distinctive spectral artifacts, partly associated with upsampling operations~\cite{frank2020frequency}. FreqNet and later frequency-domain variants extend this idea with frequency-aware representation learning, dual-frequency cues, spectral learning, and bit-plane modeling~\cite{tan2024freqnet,yan2025dffreq,karageorgiou2025spai,wang2025lota}. These methods are attractive because they move detection away from object-level shortcuts, but their signals can still be affected by compression, resizing, and post-processing.

A complementary low-level strategy anchors the decision on real-image statistics. Detecting Generated Images by Real Images, often abbreviated as LNP in benchmark reports, models noise and frequency patterns of real images for cross-generator authentication~\cite{liu2022realimages}, and its real-image-only extension moves toward one-class detection~\cite{bi2023realonly}. Beyond Generation estimates a real-image feature distribution using diffusion-model features~\cite{zhong2025beyond}, while ID-energy detection treats synthetic patterns as out-of-distribution samples relative to natural-image patterns~\cite{cai2025generalizable}. SDAIE further uses camera metadata as a self-supervised signal and can be evaluated in one-class or binary-source settings~\cite{zhong2026metadata}; related camera-pipeline work exploits demosaicing-guided color correlations~\cite{zhong2026color}. These studies support our use of low-level priors, while highlighting that camera and frequency cues should be combined with other evidence rather than used as the only basis for a fixed threshold.

\subsection{Gradient, Residual, and Reconstruction Cues}
Residual and reconstruction methods use the behavior of auxiliary models or image operators to expose traces that may not be visible in RGB space. LGrad represents an image through gradients of a pretrained model, aiming to capture more general artifact evidence~\cite{tan2023lgrad}. NPR revisits nearest-neighbor relationships induced by upsampling and models neighboring-pixel regularities shared across GAN and diffusion images~\cite{tan2024npr}. These approaches are closely related to our residual prior, but our use is deliberately constrained: residual evidence is treated as a calibrated correction signal instead of a standalone detector.

Reconstruction-based detectors are especially prominent for diffusion-generated images. DIRE uses diffusion reconstruction error as evidence of synthesis~\cite{wang2023dire}; AEROBLADE uses autoencoder reconstruction error for training-free latent-diffusion detection~\cite{ricker2024aeroblade}; LaRE\textsuperscript{2} moves reconstruction-error analysis into latent space~\cite{luo2024lare2}; FakeInversion detects unseen text-to-image models through Stable Diffusion inversion~\cite{cazenavette2024fakeinversion}; and FIRE introduces frequency-guided reconstruction error for robust diffusion-image detection~\cite{chu2025fire}. Related zero-shot work exploits generative-model bias through model-agnostic training-free detection, score-function geometry, or denoising-trajectory differences~\cite{he2024rigid,brokman2025manifold,liang2025trajectory}. These methods provide strong evidence, but their scores can depend on the chosen inverse process, autoencoder, or diffusion prior. This motivates combining residual or reconstruction cues with artifact and semantic evidence under a locked operating rule.

\subsection{Semantic and Hybrid Detectors}
Foundation models have shifted the field from detector-specific feature learning toward transferable representation spaces. CLIP provides large-scale vision-language features that transfer across visual tasks~\cite{radford2021clip}. UnivFD showed that simple nearest-neighbor or linear-probe rules in a frozen vision-language feature space can generalize across generator families better than conventional real-vs-fake classifiers trained on one fake source~\cite{ojha2023univfd}. Follow-up CLIP-based work found that few examples from one generator can provide strong out-of-distribution and robustness performance~\cite{cozzolino2024clip}, while C2P-CLIP improves CLIP-based generalization by injecting category-common prompts~\cite{tan2025c2pclip}. The strength of this paradigm is that the representation is not learned solely from one real/fake split; its limitation is that semantic transfer alone may miss low-level traces.

Recent hybrid detectors therefore combine representation-level and low-level evidence. AIDE integrates CLIP-level semantics with patch-level features and shows that visually challenging generated images remain difficult for existing detectors~\cite{yan2025aide}. CO-SPY combines semantic and pixel-level features~\cite{cheng2025cospy}; D$^3$ uses discrepancy learning to scale detector training~\cite{yang2025d3}; RINE leverages intermediate encoder-block representations~\cite{koutlis2024rine}; EFFORT decomposes CLIP features into orthogonal subspaces for generalizable detection~\cite{yan2024effort}; MCAN aggregates multiple cue experts~\cite{tan2026mcan}; and SynerDetect aligns perceptual and generative forensic representations in a unified detection space~\cite{li2026synerdetect}. This trend supports multi-evidence detection, but it also raises a protocol question: if multiple cues are fused, the fusion rule and operating threshold must be fixed without target-label feedback.

\subsection{Benchmarks and Calibration Protocols}
As generators diversify, evaluation protocol has become inseparable from detector design. GenImage provides a million-scale benchmark with cross-generator and degraded-image settings~\cite{zhu2023genimage}, and DeepfakeBench emphasizes standardized pipelines for comparing deepfake detectors~\cite{yan2023deepfakebench}. AIDE's Chameleon analysis and recent real-world robustness benchmarks show that detector behavior can change across visually challenging samples, scenario shift, degradation, and semantic diversity~\cite{yan2025aide,li2025rrdataset}. Calibration is another protocol risk: threshold choice and score calibration can substantially change reported accuracy~\cite{yang2026calibrated}. These studies indicate that high benchmark accuracy is not sufficient if the score, model variant, fusion constant, or decision threshold is selected using target information. Prior work provides strong single-signal detectors, broad benchmarks, and increasingly capable hybrid models, but the source-only integration protocol for complementary cues remains under-specified. We therefore study artifact, semantic, and residual priors under a locked fusion and threshold rule, reserving target labels only for reporting.

\section{Method}
\label{sec:method}

\subsection{Detection Setting and Source-Only Protocol}
\label{sec:problem_formulation}

We begin by defining the detection setting and the source-only rule that constrains all later design choices.

\textbf{Task definition.}
We study binary detection of AI-generated images. Given an RGB image $x \in \mathcal{X}$, the label $y \in \{0,1\}$ indicates whether the image is real ($y=0$) or generated by an AI model ($y=1$). A detector produces a scalar fake score
\begin{equation}
    f(x) \in [0,1],
\end{equation}
where larger values indicate stronger evidence that $x$ is generated. Unless stated otherwise, the final decision uses a fixed operating threshold
\begin{equation}
    \hat{y}(x) = \mathbb{I}\left[f(x) \geq 0.50\right].
\end{equation}
This fixed-threshold rule is part of the evaluated detector rather than a post-hoc choice made on a target benchmark.

\textbf{Source and target domains.}
Let $\mathcal{D}_{s}=\{(x_i^s,y_i^s)\}_{i=1}^{n_s}$ denote the source data available before deployment. In our implementation, the source data come from the ProGAN training split with balanced real and generated images. We separate source use into two auditable roles: source-fitted component training, denoted by $\mathcal{D}_{s}^{\mathrm{fit}}$ when a fit/gate split is explicitly materialized, and a recorded source-gate subset $\mathcal{D}_{s}^{\mathrm{gate}}$ used to select fusion constants from cached component scores. The recorded materialization uses seed 100 and assigns 80,000 source images to the fit role and 20,000 to the gate role, with 2,000 fit and 500 gate images for each class-label cell. The promoted runtime checkpoints are therefore described as source-fitted checkpoints, while the promoted fusion constants are source-gate selected. Let $\mathcal{D}_{t}^{g}=\{(x_i^{t,g},y_i^{t,g})\}_{i=1}^{n_g}$ denote a report-only target set associated with a generator or benchmark condition $g$. The target domain may differ from the source domain in generator architecture, image content, synthesis pipeline, resolution, and post-processing. The goal is therefore not in-domain source accuracy alone, but transfer of a locked detector from $\mathcal{D}_{s}$ to unseen target generators $\mathcal{G}_{t}$.

\textbf{Source-only selection rule.}
The central protocol constraint is that all detector choices must be made without target labels. Source labels may be used to fit detector priors, fit the CLIP linear probe, choose the residual checkpoint from source-side training records, and select fusion constants on $\mathcal{D}_{s}^{\mathrm{gate}}$. The operating threshold is fixed in advance at $0.50$ and is not tuned on target data. Target labels are used only after the detector is locked, to compute evaluation metrics. In particular, target labels must not determine prior training, epoch or checkpoint choice, fusion weights, score normalization, threshold selection, generator subset selection, or which score variant is reported. This rule treats target benchmarks as diagnostics of transfer rather than as validation sets for detector selection.

\textbf{Evaluation setting and scope.}
For each target generator or benchmark condition, we report ranking metrics such as average precision and area under the receiver operating characteristic curve, together with fixed-threshold metrics computed at $\tau=0.50$. This separation is important because a detector can rank real and generated images well while still requiring a different operating threshold under domain shift. We evaluate image-level binary authenticity from RGB pixels. The formulation does not assume access to generator internals, prompts, provenance metadata, watermarks, or camera metadata at test time, and it does not address localized manipulation, model attribution, or face presentation attack detection.

\subsection{Overview}
\label{sec:method_overview}

FreqPRISM is a source-only multi-prior detector for cross-generator fake-image detection. Its design is driven by a simple deployment constraint: after the detector is locked on the source domain, target labels must not change its component priors, fusion constants, operating threshold, or reported score variant. We therefore organize the method as a two-stage pipeline, illustrated in Fig.~\ref{fig:freqprism_overview}, that separates source-side locking from target-side inference.

In the locking stage, source-fitted artifact, semantic, and residual priors are evaluated on a recorded source-gate split to build a component-score cache. This cache is used only to select compact fusion constants under the fixed threshold $\tau=0.50$, producing the detector that will be deployed. In the inference stage, the locked detector maps each image to four scalar cues: whole-image artifact evidence, native-tile artifact evidence, CLIP semantic evidence, and nearest-neighbor residual evidence. These cues are combined in logit space with the locked constants, so target benchmarks serve only as report-time evaluations rather than validation sets.

\begin{figure*}[t]
\centering
\includegraphics[width=\textwidth]{Fig/freqprism_overview_placeholder.png}
\caption{Overview of FreqPRISM. The upper branch shows source-only locking: source-fitted priors produce a source-gate component cache, from which compact fusion constants are selected under a fixed threshold. The lower branch shows locked inference: an input image is scored by whole-image artifact, native-tile artifact, semantic, and residual priors, then fused in logit space to produce the final fake score and fixed-threshold decision.}
\label{fig:freqprism_overview}
\end{figure*}

The remainder of this section follows this pipeline. Section~\ref{sec:method_scoring} defines the multi-prior scoring abstraction and the source-only locking routine. Section~\ref{sec:method_artifact} describes the artifact and tile priors, Sections~\ref{sec:method_semantic} and~\ref{sec:method_residual} describe the semantic and residual priors, and Section~\ref{sec:method_fusion} defines the final fusion rule and locked inference decision. Table~\ref{tab:component_priors} summarizes the component-level implementation choices that are fixed before inference.

\begin{table*}[t]
\caption{Source-Fitted Priors and Locked Score Construction in FreqPRISM}
\label{tab:component_priors}
\centering
\footnotesize
\renewcommand{\arraystretch}{1.12}
\setlength{\tabcolsep}{3pt}
\begin{tabular}{p{0.14\textwidth}p{0.18\textwidth}p{0.26\textwidth}p{0.17\textwidth}p{0.16\textwidth}}
\toprule
Score & Input view & Source-fitted model & Locked preprocessing & Output role \\
\midrule
$W(x)$ & Whole-image artifact views & 849-dimensional explicit artifact features with codec histogram-gradient boosting and chroma logistic experts; artifact payload coefficient $\alpha_A=-0.40$ & Bicubic $256\times256$ clean, JPEG quality 35 and 50, half-scale resize, and Gaussian blur radius 1 views & Whole-image low-level authenticity score aggregated by mean logit \\
$T(x)$ & Native image tiles & Same clean artifact branch as $W(x)$ & $256\times256$ tiles from a $3\times3$ native grid, with edge padding for boundary crops and top-1 aggregation & Positive local artifact increment relative to $W(x)$ \\
$S(x)$ & CLIP clean view & Frozen OpenAI CLIP ViT-L/14 image encoder with a source-fitted standardized balanced logistic-regression probe, $C=1.0$ & Clean $256\times256$ image view with L2-normalized CLIP features & Signed semantic evidence in logit space \\
$R(x)$ & Nearest-neighbor residual tensor & ResNet-50 binary residual prior trained on source data with binary cross-entropy for two epochs; runtime loads checkpoint 1 & RGB conversion, even-size crop, tensor conversion, ImageNet normalization, and nearest-neighbor down-up residual & Low-weight residual correction with coefficient $\gamma$ \\
\bottomrule
\end{tabular}
\end{table*}

\subsection{Source-Only Multi-Prior Scoring}
\label{sec:method_scoring}

Under the source-only protocol, the detector must convert heterogeneous cues into one deployable score without using target labels for calibration or model selection. FreqPRISM therefore uses a fixed multi-prior scoring pipeline. For an input image $x$, four source-fitted components produce scalar scores: a whole-image artifact score $W(x)$ from explicit low-level artifact features, a native-tile artifact score $T(x)$ from local high-resolution evidence, a semantic score $S(x)$ from a frozen CLIP image encoder with a source-fitted linear probe, and a residual score $R(x)$ from a nearest-neighbor down-up reconstruction residual branch. We use these quantities as bounded component scores and map them to clipped logits only inside the fusion rule. The source-gate stress constraints limit score drift and fixed-threshold flips after fusion, rather than claiming that the heterogeneous component probabilities are perfectly calibrated across domains. The resulting locked fake score is denoted by $f(x)$.

The four scores are not treated as interchangeable ensemble members. The whole-image artifact prior anchors the decision with low-level authenticity evidence; the native-tile pathway contributes only when local artifact evidence is stronger than the whole-image estimate; the semantic prior supplies signed representation-level evidence from a pretrained vision-language feature space; and the residual prior adds a controlled correction for nearest-neighbor down-up reconstruction traces. This design keeps the detector interpretable at the component level while preserving the protocol constraint that all fusion choices are fixed before target evaluation.

Algorithm~\ref{alg:source_locking} summarizes how the fusion constants are selected before target evaluation. The algorithm operates only on the source-gate cache and the fixed threshold. It does not inspect target scores, target labels, or per-generator target summaries.

\begin{algorithm}[t]
\caption{Source-only fusion calibration}
\label{alg:source_locking}
\begin{algorithmic}[1]
\REQUIRE Source-gate component cache $\mathcal{C}_{s}^{\mathrm{gate}}$, anchor constants $\theta_0$, scale grid $\mathcal{G}$, source constraints $\Omega_s$, fixed threshold $\tau=0.50$
\ENSURE Locked fusion constants $\theta^\star$ and threshold $\tau$
\STATE Load $\mathcal{C}_{s}^{\mathrm{gate}}=\{(W_i,T_i,S_i,R_i,m_i,y_i)\}$ from source-gate component scores.
\STATE $\mathcal{A}\leftarrow\emptyset$
\FOR{each $(s_T,s_{S+},s_{S-},s_R)\in\mathcal{G}^4$}
    \STATE $\theta\leftarrow\textsc{FoldScales}(\theta_0,s_T,s_{S+},s_{S-},s_R)$
    \STATE $f_i(\theta)\leftarrow\textsc{Fuse}(W_i,T_i,S_i,R_i,m_i;\theta)$ for each gate sample.
    \STATE Compute source ranking, threshold, drift, flip-rate, and source log-loss metrics.
    \IF{all constraints in $\Omega_s$ are satisfied}
        \STATE Add $\theta$ and its source metrics to $\mathcal{A}$.
    \ENDIF
\ENDFOR
\STATE Select $\theta^\star\in\mathcal{A}$ with minimum fake-side source log loss, breaking ties by anchor distance, real-side log loss, source log loss, and variant identifier.
\RETURN $\theta^\star,\tau$
\end{algorithmic}
\end{algorithm}

\subsection{Artifact Prior and Native-Tile Evidence}
\label{sec:method_artifact}

The artifact prior converts explicit low-level image statistics into an image-level fake score. For an input image, we first construct a $256 \times 256$ artifact view and extract a feature vector $\phi(x)$. The feature extractor includes statistics from RGB-light, residual, high-pass, local-DCT, and global-FFT views, together with rich spatial statistics, texture artifacts, recompression stability, residual co-occurrence and spectrum, residual-tail shape, chroma-luma coupling, patch-spectrum heterogeneity, and codec-block statistics. In the current implementation, $\phi(x)$ has 849 dimensions. Two source-fitted experts are trained from source labels: a codec expert $C(\cdot)$, implemented as a histogram gradient boosting classifier on codec-block features, and a chroma expert $Q(\cdot)$, implemented as a standardized logistic regression probe on fine chroma-luma coupling features. Both experts output fake probabilities. The artifact training record uses expanded source variants consisting of clean images, two JPEG-quality-50 copies, half-scale resize, and Gaussian blur radius 1, with no target labels.

For a probability $p$, let
\begin{equation}
    \ell(p)=\log \frac{p}{1-p}, \qquad
    \sigma(z)=\frac{1}{1+\exp(-z)},
\end{equation}
with probabilities clipped to $[10^{-6},1-10^{-6}]$ in implementation. For each artifact evaluation variant $v \in \mathcal{V}_{A}$, where $\mathcal{V}_{A}$ contains bicubic-resized clean, JPEG quality 35, JPEG quality 50, half-scale resized, and Gaussian-blurred $256 \times 256$ views, the artifact branch computes
\begin{equation}
    A_v(x)=\sigma\left(\ell(C(\phi(v(x))))+\alpha_A \ell(Q(\phi(v(x))))\right),
\end{equation}
where $\alpha_A=-0.40$ is fixed in the artifact-prior payload. The whole-image artifact score is the mean-logit aggregation
\begin{equation}
    W(x)=\sigma\left(\frac{1}{|\mathcal{V}_{A}|}\sum_{v\in\mathcal{V}_{A}}\ell(A_v(x))\right).
\end{equation}
This aggregation uses perturbation views to stabilize the low-level score against common compression and resizing changes. The artifact experts are fitted from the expanded source variants recorded in the training protocol and do not use target images or target labels.

Whole-image resizing can dilute local high-resolution artifacts. FreqPRISM therefore computes a native-tile artifact score in parallel. If both image sides are no larger than the tile size, the tile branch uses one RGB crop covering the image and pads it by edge replication to $256 \times 256$ when necessary. Otherwise, the image is cropped into native-resolution $256 \times 256$ tiles from a $3 \times 3$ grid whose start positions include both image boundaries on each axis. Boundary tiles smaller than $256 \times 256$ are padded by edge replication before scoring. Each tile is scored by the clean artifact branch, and the image-level tile score is
\begin{equation}
    T(x)=\max_{j} A_{\mathrm{clean}}(t_j),
\end{equation}
where $\{t_j\}$ are the extracted native tiles. Tile evidence enters the fused detector only as a positive logit increment over the whole-image score:
\begin{equation}
    \Delta_T(x)=\max\left(0,\ell(T(x))-\ell(W(x))\right).
\end{equation}
This design lets a strong local artifact increase fake evidence while preventing low-scoring tiles from suppressing a reliable whole-image artifact score.

\subsection{Semantic Prior}
\label{sec:method_semantic}

The semantic prior supplies representation-level evidence that is not learned from low-level artifact statistics alone. We use a frozen CLIP ViT-L/14 image encoder~\cite{radford2021clip}. The CLIP encoder is not fine-tuned. Given an image $x$, the frozen encoder produces an image representation
\begin{equation}
    z(x)=E_{\mathrm{CLIP}}(x).
\end{equation}
A source-only logistic regression probe $g_S$ is then fitted using CLIP features and source labels. The semantic fake score is
\begin{equation}
    S(x)=\sigma(g_S(z(x))).
\end{equation}
The probe uses L2-normalized CLIP features, a standardization layer, and a balanced logistic regression classifier with fixed regularization $C=1.0$ on the clean image view. The training script records source-side diagnostic metrics for the probe, but these diagnostics are not used to choose the CLIP backbone, regularization, score variant, fusion constants, or threshold. Target labels are not used to choose the CLIP backbone, the probe regularization, the score variant, or the final semantic contribution. In the fusion rule, $S(x)$ is treated as signed evidence: positive semantic logits support the fake class, while negative semantic logits support the real class. Separate coefficients are used for positive and negative semantic evidence because their calibration roles differ under domain shift.

\subsection{Residual Prior}
\label{sec:method_residual}

The residual prior adds a correction signal based on nearest-neighbor down-up traces. This branch follows the observation that neighboring-pixel and upsampling residuals expose artifacts that may not be captured by RGB classifiers alone~\cite{tan2024npr}. Let $\tilde{x}$ denote the residual-branch tensor after RGB conversion, even-size cropping, tensor conversion, and ImageNet normalization. During training, $\tilde{x}$ is resized to $256 \times 256$; in the locked inference configuration, the residual branch uses the native even-cropped resolution. Let $D_{\mathrm{nn}}$ and $U_{\mathrm{nn}}$ denote nearest-neighbor downsampling and upsampling by a factor of two. The residual input is
\begin{equation}
    \rho(x)=\tilde{x}-U_{\mathrm{nn}}(D_{\mathrm{nn}}(\tilde{x})).
\end{equation}
A ResNet-50 binary classifier $g_R$ maps this residual tensor to a fake logit, and the residual score is
\begin{equation}
    R(x)=\sigma(g_R(\rho(x))).
\end{equation}
The residual model is trained on source images for two epochs with binary cross-entropy, and the locked runtime configuration uses the epoch-1 checkpoint. This checkpoint choice is part of the source-side configuration record and is fixed before any target report is generated. In the final detector, the residual branch is not used as a standalone classifier. It enters only through the low-weight logit correction in Section~\ref{sec:method_fusion}, which limits the influence of residual scores that may vary with the reconstruction operator and preprocessing resolution.

\subsection{Source-Only Fusion and Inference}
\label{sec:method_fusion}

FreqPRISM fuses the four prior scores in logit space. Let $m(x)$ indicate whether $x$ is high resolution, defined by maximum side length greater than 960 pixels, and let $\Delta_T(x)$ be the positive tile increment defined above. The semantic coefficient for positive evidence is
\begin{equation}
    \alpha_{+}(x)=
    \begin{cases}
    \alpha_{\mathrm{high}+}, & m(x)=1,\\
    \alpha_{\mathrm{low}+}, & m(x)=0,
    \end{cases}
\end{equation}
and the semantic coefficient for negative evidence is
\begin{equation}
    \alpha_{-}(x)=
    \begin{cases}
    \alpha_{\mathrm{high}-}, & m(x)=1 \ \mathrm{and}\ \Delta_T(x)>0,\\
    \alpha_{\mathrm{high}-}^{\mathrm{guard}}, & m(x)=1 \ \mathrm{and}\ \Delta_T(x)\leq 0,\\
    \alpha_{\mathrm{low}-}, & m(x)=0.
    \end{cases}
\end{equation}
The semantic logit contribution is
\begin{equation}
    \Psi_S(x)=\alpha_{+}(x)\max(0,\ell(S(x)))+
    \alpha_{-}(x)\min(0,\ell(S(x))).
\end{equation}
The base score before residual correction is
\begin{equation}
    B(x)=\sigma\left(\ell(W(x))+\beta \Delta_T(x)+\Psi_S(x)\right),
\end{equation}
and the final fake score is
\begin{equation}
    f(x)=\sigma\left(\ell(B(x))+\gamma \ell(R(x))\right).
\end{equation}

The fusion rule separates its functional form from its locked numerical instantiation. We select the constants once from source-gate component scores and source labels only, using the search in Algorithm~\ref{alg:source_locking}. The source-gate search uses the compact scale grid $\{0.75,1.00,1.25,1.50,1.75,2.00\}$ for the tile, positive semantic, negative semantic, and residual terms. A candidate is accepted only if it satisfies the source BA, AP, AUC, fixed-threshold flip-rate, mean-drift, and real-side log-loss constraints defined in the source-locking record. This search accepts 231 of 1,296 compact-scale candidates and selects the candidate with tile scale 1.25, positive semantic scale 2.00, negative semantic scale 1.25, and residual scale 1.75 relative to the source-gamma anchor. After folding these scales into the runtime configuration, the selected configuration uses $\beta=0.25$ for the tile increment and $\gamma=0.21$ for the residual correction. For semantic evidence, it uses $\alpha_{\mathrm{low}+}=0.30$ and $\alpha_{\mathrm{low}-}=0.1875$ on low-resolution images, and $\alpha_{\mathrm{high}+}=0.40$, $\alpha_{\mathrm{high}-}=0.00$, and $\alpha_{\mathrm{high}-}^{\mathrm{guard}}=0.25$ on high-resolution images. These values are fixed before any target report is generated.

The final decision is
\begin{equation}
    \hat{y}(x)=\mathbb{I}\left[f(x)\geq 0.50\right],
\end{equation}
where the threshold is not tuned on target generators. Thus the target benchmarks evaluate a locked detector, rather than selecting fusion parameters, score variants, or an operating point. Algorithm~\ref{alg:locked_inference} summarizes the complete locked inference procedure for a test image.

\begin{algorithm}[t]
\caption{Locked multi-prior scoring}
\label{alg:locked_inference}
\begin{algorithmic}[1]
\REQUIRE Image $x$, locked priors $\mathcal{P}$, locked constants $\theta^\star$, fixed threshold $\tau=0.50$
\ENSURE Fake score $f(x)$ and binary decision $\hat{y}(x)$
\STATE Compute whole-image artifact score $W(x)$ by mean-logit aggregation over $\mathcal{V}_A$.
\STATE Extract native $256\times256$ tiles from $x$; use edge padding for boundary tiles when needed.
\STATE Score each tile with the clean artifact branch and set $T(x)\leftarrow\max_j A_{\mathrm{clean}}(t_j)$.
\STATE Compute semantic score $S(x)$ using frozen CLIP features and the source-fitted linear probe.
\STATE Compute residual score $R(x)$ from nearest-neighbor down-up residuals.
\STATE $\Delta_T(x)\leftarrow\max(0,\ell(T(x))-\ell(W(x)))$
\STATE $m(x)\leftarrow\mathbb{I}[\max\mathrm{side}(x)>960]$
\STATE Set $\alpha_+(x)$ from $m(x)$ and set $\alpha_-(x)$ using the high-resolution guard.
\STATE $u_S\leftarrow\max(0,\ell(S(x)))$, \quad $v_S\leftarrow\min(0,\ell(S(x)))$
\STATE $\Psi_S(x)\leftarrow\alpha_+(x)u_S+\alpha_-(x)v_S$
\STATE $B(x)\leftarrow\sigma(\ell(W(x))+\beta\Delta_T(x)+\Psi_S(x))$
\STATE $f(x)\leftarrow\sigma(\ell(B(x))+\gamma\ell(R(x)))$
\STATE $\hat{y}(x)\leftarrow\mathbb{I}[f(x)\geq\tau]$
\RETURN $f(x),\hat{y}(x)$
\end{algorithmic}
\end{algorithm}

\section{Experiments}
\label{sec:experiments}
We evaluate FreqPRISM under two source protocols. The primary protocol follows the common ProGAN-to-unseen-generator setting: the detector is fitted and locked with ProGAN source data, then reported on AIGCDetectBenchmark~\cite{zhong2023patchcraft} and UniversalFakeDetect. The second protocol follows the GenImage setting, where the source generator is Stable Diffusion v1.4 and the detector is evaluated across the GenImage validation generators. In both protocols, target labels are reserved for final reporting and are not used to select components, fusion constants, calibration parameters, or thresholds.

\subsection{Experimental Setup}
\label{sec:experimental_setup}

\textbf{Datasets.}
We evaluate cross-generator generalization on three public benchmarks. AIGCDetectBenchmark~\cite{zhong2023patchcraft} follows the ProGAN-to-unseen-generator protocol and contains GAN, diffusion, and commercial-generator test subsets; its generator composition is summarized in Table~\ref{tab:aigcdetectbenchmark}. UniversalFakeDetect~\cite{ojha2023univfd} is used as an additional ProGAN-source transfer benchmark with the same locked detector. GenImage~\cite{zhu2023genimage} provides an SD v1.4-source protocol and evaluates transfer to eight generator families, covering both diffusion models and BigGAN.

\begin{table}[t]
\caption{The description of generative models. SD and WFIR denote Stable Diffusion and whichfaceisreal, respectively.}
\label{tab:aigcdetectbenchmark}
\centering
\scriptsize
\renewcommand{\arraystretch}{1.04}
\setlength{\tabcolsep}{6pt}
\begin{tabular}{@{}cccc@{}}
\toprule
\multirow{2}{*}{Generator} & Image & Image & Photographic \\
 & Size & Number & Source \\
\midrule
ProGAN~\cite{karras2018progan} & 256$\times$256 & 8.0k & LSUN~\cite{yu2015lsun} \\
StyleGAN~\cite{karras2019stylegan} & 256$\times$256 & 12.0k & LSUN \\
BigGAN~\cite{brock2019biggan} & 256$\times$256 & 4.0k & ImageNet~\cite{deng2009imagenet} \\
CycleGAN~\cite{zhu2017cyclegan} & 256$\times$256 & 2.6k & ImageNet \\
StarGAN~\cite{choi2018stargan} & 256$\times$256 & 4.0k & CelebA~\cite{liu2015celeba} \\
GauGAN~\cite{park2019spade} & 256$\times$256 & 10.0k & COCO~\cite{lin2014coco} \\
StyleGAN2~\cite{karras2020stylegan2} & 256$\times$256 & 15.9k & LSUN \\
WFIR~\cite{whichfaceisreal2023} & 1024$\times$1024 & 2.0k & FFHQ~\cite{karras2019stylegan} \\
ADM~\cite{dhariwal2021diffusion} & 256$\times$256 & 12.0k & ImageNet \\
Glide~\cite{nichol2022glide} & 256$\times$256 & 12.0k & ImageNet \\
Midjourney~\cite{midjourney2023} & 1024$\times$1024 & 12.0k & ImageNet \\
SD v1.4~\cite{rombach2022ldm} & 512$\times$512 & 12.0k & ImageNet \\
SD v1.5~\cite{rombach2022ldm} & 512$\times$512 & 16.0k & ImageNet \\
VQDM~\cite{gu2022vqdm} & 256$\times$256 & 12.0k & ImageNet \\
WUKONG~\cite{wukong2023} & 512$\times$512 & 12.0k & ImageNet \\
DALL-E 2~\cite{ramesh2022dalle2} & 256$\times$256 & 2.0k & ImageNet \\
SDXL~\cite{podell2024sdxl} & 1024$\times$1024 & 4.0k & COCO \\
\bottomrule
\end{tabular}
\end{table}

\textbf{Training and locking protocols.}
For AIGCDetectBenchmark and UniversalFakeDetect, we follow the ProGAN-source setting but use a balanced 100k-image source subset, consisting of 50k real images and 50k ProGAN images. This source budget is smaller than the commonly used CNNSpot ProGAN training set for AIGCDetectBenchmark, which contains 360k LSUN real images and 360k ProGAN images~\cite{wang2020cnnspot,zhong2023patchcraft}. All components are fitted from this source subset, and the fusion constants are locked on source-only evidence before target evaluation. The same locked detector is then evaluated on AIGCDetectBenchmark and UniversalFakeDetect. For GenImage, the components are fitted with the SD v1.4 source data and the operating score is calibrated using only source-calibration images. In all protocols, target labels are used only for final reporting; they are never used to choose a component variant, fusion constant, calibration bias, or decision threshold.

\textbf{Implementation details.}
All reports use the FreqPRISM scoring pipeline described in Section~\ref{sec:method}: whole-image artifact scoring, native-tile artifact scoring, frozen-CLIP semantic scoring, and nearest-neighbor residual scoring followed by locked logit fusion. The ProGAN-source reports use the folded constants in Section~\ref{sec:method_fusion} and the fixed threshold $\tau=0.50$. The GenImage report uses the same fusion form with source-only calibration from the SD v1.4 source data before applying the same threshold.

\textbf{Metrics.}
Following common reports on AIGCDetectBenchmark, UniversalFakeDetect, and GenImage, we use classification accuracy (Acc) and average precision (AP) as the primary comparable metrics~\cite{ojha2023univfd,zhong2023patchcraft,zhu2023genimage,yan2025aide}. We additionally report the area under the ROC curve (AUC) and the real/fake accuracy split. Acc, Real, and Fake are threshold metrics computed at $\tau=0.50$ after the source-side locking or calibration step, whereas AP and AUC are threshold-free ranking metrics. AUC complements AP by summarizing ROC ordering across operating points, and the Real/Fake split exposes asymmetric calibration errors that can be hidden by mean Acc. These added diagnostics are reported after locking and are not used for model selection.

\subsection{Benchmark-Level Reports}
\label{sec:benchmark_reports}

\begin{table}[t]
\caption{Cross-Benchmark Evaluation Under Locked Source Protocols. Acc, AP, AUC, Real, and Fake are averaged over target generators within each benchmark.}
\label{tab:benchmark_summary}
\centering
\scriptsize
\renewcommand{\arraystretch}{1.08}
\setlength{\tabcolsep}{2.8pt}
\begin{tabular}{@{}lcrrrrr@{}}
\toprule
Benchmark & Source & Target & Mean & Mean & Mean & Mean \\
 & protocol & generators & Acc & AP & AUC & Real/Fake \\
\midrule
AIGCDetectBenchmark & ProGAN & 17 & 94.82 & 99.34 & 99.30 & 94.69/94.95 \\
UniversalFakeDetect & ProGAN & 21 & 89.39 & 99.24 & 99.19 & 83.36/95.41 \\
GenImage & SD v1.4 & 8 & 93.84 & 98.50 & 98.56 & 94.91/92.76 \\
\bottomrule
\end{tabular}
\end{table}

Table~\ref{tab:benchmark_summary} gives the compact report across the three evaluated benchmarks, with metrics averaged over target generators. On AIGCDetectBenchmark, the locked ProGAN-source detector attains 94.82 Acc and 99.34 AP across 17 generators. On UniversalFakeDetect, the same ProGAN-source detector keeps high ranking quality but shows a lower real-side fixed-threshold accuracy, which lowers mean Acc despite 99.24 AP. On GenImage, the SD v1.4-source protocol gives 93.84 Acc and 98.50 AP across eight generator families.

These results motivate two reporting choices. First, AP and AUC are not sufficient for this detector, because the fixed decision threshold exposes benchmark-specific shifts. Second, the source protocol must be stated with each benchmark. AIGCDetectBenchmark and UniversalFakeDetect evaluate the same ProGAN-source locked detector, whereas GenImage evaluates the same FreqPRISM design with a GenImage SD v1.4 source and source-only probability calibration.

\subsection{Cross-Generator Behavior}
\label{sec:cross_generator_behavior}

Table~\ref{tab:aigcdetect_comparison} reports generator-level transfer on AIGCDetectBenchmark~\cite{zhong2023patchcraft}. The main text uses Acc because recent AIGCDetectBenchmark reports provide complete 17-generator Acc more consistently than per-generator AP. The full AP counterpart is provided after the references in the reproducibility checklist (Table~\ref{tab:aigcdetect_ap_comparison}) using the same layout. The best and second-best reported values in each column are marked in bold and underlined, respectively.

\begin{table*}[!t]
\caption{Comparison on AIGCDetectBenchmark. Each entry reports Acc (\%) for detecting real and fake images from the generator named in the column. DIRE-G denotes the ProGAN-source DIRE variant, whereas DIRE-D follows the ADM-source DIRE setup; SDAIE is the one-class variant, whereas SDAIE$^\dagger$ is the binary ProGAN-source variant. The best and second-best values in each column are marked in bold and underline, respectively.}
\label{tab:aigcdetect_comparison}
\centering
\scriptsize
\providecommand{\rotgen}[1]{\makebox[0pt][l]{\rotatebox{60}{\textit{#1}}}}
\renewcommand{\arraystretch}{1.03}
\setlength{\tabcolsep}{2.0pt}
\begin{tabularx}{\textwidth}{@{}l>{\centering\arraybackslash}p{0.072\textwidth}*{18}{Y}@{}}
\toprule
Method & Ref. & \rotgen{ProGAN} & \rotgen{StyleGAN} & \rotgen{BigGAN} & \rotgen{CycleGAN} & \rotgen{StarGAN} & \rotgen{GauGAN} & \rotgen{StyleGAN2} & \rotgen{WFIR} & \rotgen{ADM} & \rotgen{GLIDE} & \rotgen{Midjourney} & \rotgen{SD v1.4} & \rotgen{SD v1.5} & \rotgen{VQDM} & \rotgen{Wukong} & \rotgen{DALL-E 2} & \rotgen{SDXL} & \rotgen{Mean} \\
\midrule
CNNSpot & CVPR~2020 & \textbf{100.0} & 90.2 & 71.2 & 87.6 & 94.6 & 81.4 & 86.9 & 91.7 & 60.4 & 58.1 & 51.4 & 50.6 & 50.5 & 56.5 & 51.0 & 50.5 & 53.0 & 69.73 \\
FreDect & ICML~2020 & 99.4 & 78.0 & 82.0 & 78.8 & 94.6 & 80.6 & 66.2 & 50.8 & 63.4 & 54.1 & 45.9 & 38.8 & 39.2 & 77.8 & 40.3 & 34.7 & 51.2 & 63.28 \\
Fusing & ICIP~2022 & \textbf{100.0} & 85.2 & 77.4 & 87.0 & 97.0 & 77.0 & 83.3 & 66.8 & 49.0 & 57.2 & 52.2 & 51.0 & 51.4 & 55.1 & 51.7 & 52.8 & 55.6 & 67.63 \\
GramNet & CVPR~2020 & \textbf{100.0} & 87.1 & 67.3 & 86.1 & 95.1 & 69.4 & 87.3 & 86.8 & 58.6 & 54.5 & 50.0 & 51.7 & 52.2 & 52.9 & 50.8 & 49.3 & 64.5 & 68.43 \\
LNP & ECCV~2022 & \textbf{100.0} & 92.6 & 88.4 & 79.1 & \textbf{100.0} & 79.2 & 93.8 & 50.0 & 83.9 & 83.5 & 69.6 & 89.3 & 88.8 & 85.0 & 86.4 & 92.5 & 87.8 & 85.28 \\
LGrad & CVPR~2023 & 99.8 & 91.1 & 85.6 & 86.9 & 99.3 & 78.5 & 85.3 & 55.7 & 67.2 & 66.1 & 65.4 & 63.0 & 63.7 & 73.0 & 59.6 & 65.5 & 71.3 & 75.11 \\
DIRE-G & ICCV~2023 & 95.2 & 83.0 & 70.1 & 74.2 & 95.5 & 67.8 & 75.3 & 58.1 & 75.8 & 71.8 & 58.0 & 49.7 & 49.8 & 53.7 & 54.5 & 66.5 & 55.4 & 67.90 \\
DIRE-D & ICCV~2023 & 52.8 & 51.3 & 49.7 & 49.6 & 46.7 & 51.2 & 51.7 & 53.3 & \textbf{98.3} & 92.4 & 89.5 & 91.2 & 91.6 & 91.9 & 90.9 & 92.5 & 91.3 & 72.70 \\
UnivFD & CVPR~2023 & 99.8 & 84.9 & 95.1 & 98.3 & 95.8 & \underline{99.5} & 75.0 & 86.9 & 66.9 & 62.5 & 56.1 & 63.7 & 63.5 & 85.3 & 70.9 & 50.8 & 50.7 & 76.80 \\
PatchCraft & arXiv~2023 & \textbf{100.0} & 92.8 & 95.8 & 70.2 & \textbf{100.0} & 71.6 & 89.6 & 85.8 & 82.2 & 83.8 & 90.1 & \underline{95.4} & \underline{95.3} & 88.9 & 91.1 & 96.6 & \textbf{98.4} & 89.85 \\
NPR & CVPR~2024 & 99.8 & 97.7 & 84.4 & 96.1 & 99.4 & 82.5 & 98.4 & 65.8 & 69.7 & 78.4 & 77.9 & 78.6 & 78.9 & 78.1 & 76.1 & 64.9 & 81.3 & 82.81 \\
AIDE & ICLR~2025 & \textbf{100.0} & \underline{99.6} & 84.0 & 98.5 & \underline{99.9} & 73.3 & 98.0 & \underline{94.2} & 93.4 & \underline{95.1} & 77.2 & 93.0 & 92.9 & \textbf{95.2} & \underline{93.6} & 96.6 & 92.2 & 92.73 \\
SDAIE & TPAMI~2026 & 78.6 & 84.8 & 57.9 & 66.2 & 74.5 & 55.3 & 85.0 & 71.5 & 88.9 & 93.2 & 85.6 & 89.7 & 89.7 & 90.8 & 87.5 & 89.0 & 89.4 & 81.02 \\
SDAIE$^\dagger$ & TPAMI~2026 & \textbf{100.0} & \textbf{99.8} & 90.5 & 91.9 & \textbf{100.0} & 84.0 & \textbf{99.2} & 84.5 & 92.9 & 93.1 & \underline{95.4} & \textbf{95.8} & \textbf{95.7} & 91.5 & \textbf{95.0} & \underline{97.1} & \underline{96.6} & \underline{94.30} \\
FatFormer & CVPR~2024 & \underline{99.9} & 98.4 & 99.5 & \textbf{99.8} & 99.1 & 98.4 & 97.5 & 87.2 & 85.7 & 90.9 & 71.0 & 68.7 & 68.0 & 85.5 & 82.3 & 74.2 & 70.5 & 86.86 \\
RINE & ECCV~2024 & \textbf{100.0} & 97.7 & \underline{99.6} & \underline{99.3} & 99.6 & \textbf{99.8} & 96.9 & \textbf{98.5} & 71.9 & 81.7 & 66.0 & 81.4 & 81.9 & 87.1 & 82.3 & 56.9 & 74.7 & 86.78 \\
C2P-CLIP & AAAI~2025 & \textbf{100.0} & 98.6 & \textbf{99.9} & 98.4 & 98.0 & 98.7 & 97.1 & 92.3 & 70.9 & 86.8 & 62.2 & 80.5 & 80.3 & 86.5 & 79.4 & 67.1 & 69.6 & 86.25 \\
D$^3$ & CVPR~2025 & \underline{99.9} & 90.8 & 97.8 & 98.3 & 97.6 & \textbf{99.8} & 88.9 & 84.9 & 77.4 & 74.2 & 70.4 & 73.7 & 73.8 & 89.7 & 72.3 & 70.1 & 71.1 & 84.17 \\
FIRE & CVPR~2025 & 96.9 & 87.3 & 83.4 & 86.1 & 95.1 & 79.6 & 85.2 & 85.4 & 85.7 & 84.4 & 77.9 & 85.2 & 85.2 & 79.8 & 84.8 & 83.5 & 77.6 & 84.88 \\
CO-SPY & CVPR~2025 & \underline{99.9} & 96.3 & 92.0 & 98.0 & 96.1 & 90.9 & 97.9 & 71.7 & 73.3 & 88.8 & 88.7 & 86.0 & 86.4 & 82.4 & 78.8 & 77.1 & 87.5 & 87.73 \\
\midrule
\rowcolor[gray]{0.88}
Ours & -- & \textbf{100.0} & 99.4 & 89.3 & 96.6 & \underline{99.9} & 82.8 & \underline{98.9} & 90.2 & \underline{96.3} & \textbf{98.4} & \textbf{97.2} & 92.6 & 92.9 & \underline{93.4} & 87.6 & \textbf{98.2} & \textbf{98.4} & \textbf{94.82} \\
\bottomrule
\end{tabularx}
\end{table*}

% \clearpage

% Reserve Table V for the AP reproducibility table placed after the references.
\addtocounter{table}{1}

In Table~\ref{tab:aigcdetect_comparison}, Ours gives the highest mean accuracy among the listed baselines while using the locked ProGAN-source protocol. Table~\ref{tab:aigcdetect_ap_comparison} complements this result by separating threshold-free ranking quality from fixed-threshold behavior: several methods retain high AP on individual generators even when Acc drops, confirming that both metrics are needed for cross-generator reporting.

\begin{table*}[!t]
\caption{Comparison on GenImage. Each entry reports Acc/AP (\%) for detecting real and fake images from the generator named in the column. The best and second-best values in each column are marked in bold and underline, respectively.}
\label{tab:genimage_report}
\centering
\scriptsize
\providecommand{\pairslash}[2]{#1/#2}
\renewcommand{\arraystretch}{1.08}
\setlength{\tabcolsep}{3.2pt}
\resizebox{\textwidth}{!}{%
\begin{tabular}{@{}l c *{9}{c}@{}}
\toprule
Method & Ref. & BigGAN & Midjourney & SD v1.4 & SD v1.5 & ADM & GLIDE & Wukong & VQDM & Mean \\
\midrule
LGrad & CVPR~2023 & \pairslash{40.69}{39.36} & \pairslash{73.74}{77.63} & \pairslash{76.28}{79.15} & \pairslash{77.12}{80.27} & \pairslash{51.96}{51.38} & \pairslash{49.92}{50.57} & \pairslash{73.18}{75.61} & \pairslash{52.83}{51.94} & \pairslash{61.97}{63.24} \\
DIRE & ICCV~2023 & \pairslash{50.66}{54.58} & \pairslash{77.43}{86.49} & \pairslash{94.42}{99.09} & \pairslash{94.13}{99.09} & \pairslash{72.08}{82.18} & \pairslash{71.23}{81.65} & \pairslash{93.99}{98.85} & \pairslash{55.83}{62.69} & \pairslash{76.22}{83.08} \\
UnivFD & CVPR~2023 & \pairslash{81.21}{96.13} & \pairslash{72.68}{92.47} & \pairslash{85.47}{97.08} & \pairslash{85.11}{97.09} & \pairslash{55.34}{67.91} & \pairslash{77.43}{93.66} & \pairslash{76.48}{93.63} & \pairslash{57.24}{78.48} & \pairslash{73.87}{89.55} \\
NPR & CVPR~2024 & \pairslash{81.84}{88.64} & \pairslash{\underline{89.84}}{97.61} & \pairslash{90.74}{98.77} & \pairslash{90.74}{98.82} & \pairslash{\underline{84.63}}{93.16} & \pairslash{90.31}{98.34} & \pairslash{90.67}{98.06} & \pairslash{87.04}{93.05} & \pairslash{88.23}{95.81} \\
FreqNet & AAAI~2024 & \pairslash{90.58}{94.81} & \pairslash{69.72}{78.96} & \pairslash{64.27}{74.35} & \pairslash{64.84}{75.62} & \pairslash{83.31}{91.47} & \pairslash{81.68}{88.84} & \pairslash{57.79}{66.72} & \pairslash{81.76}{89.63} & \pairslash{74.24}{82.55} \\
FatFormer & CVPR~2024 & \pairslash{\textbf{96.67}}{\underline{99.47}} & \pairslash{56.06}{62.65} & \pairslash{67.66}{81.18} & \pairslash{68.09}{80.98} & \pairslash{78.33}{91.74} & \pairslash{87.92}{95.81} & \pairslash{72.92}{85.81} & \pairslash{86.85}{96.96} & \pairslash{76.81}{86.83} \\
LaRE2 & CVPR~2024 & \pairslash{80.47}{97.40} & \pairslash{69.85}{98.08} & \pairslash{98.24}{99.88} & \pairslash{97.93}{99.80} & \pairslash{71.73}{94.87} & \pairslash{80.26}{95.14} & \pairslash{\underline{98.89}}{99.86} & \pairslash{87.48}{97.55} & \pairslash{85.61}{97.82} \\
AIDE & ICLR~2025 & \pairslash{66.83}{93.48} & \pairslash{79.38}{97.95} & \pairslash{99.75}{\underline{99.95}} & \pairslash{\underline{99.75}}{\underline{99.98}} & \pairslash{78.49}{94.58} & \pairslash{91.77}{99.08} & \pairslash{98.85}{\underline{99.97}} & \pairslash{80.24}{97.07} & \pairslash{86.88}{97.76} \\
FerretNet & NeurIPS~2025 & \pairslash{\underline{91.32}}{\textbf{99.58}} & \pairslash{78.00}{97.60} & \pairslash{\textbf{100.00}}{\textbf{100.00}} & \pairslash{\textbf{99.80}}{\textbf{100.00}} & \pairslash{76.48}{\underline{96.60}} & \pairslash{\textbf{98.56}}{\underline{99.80}} & \pairslash{\textbf{99.58}}{\textbf{99.99}} & \pairslash{81.62}{\underline{98.32}} & \pairslash{90.67}{\textbf{98.99}} \\
EFFORT & ICML~2025 & \pairslash{77.60}{98.38} & \pairslash{82.40}{\underline{98.96}} & \pairslash{\underline{99.80}}{99.82} & \pairslash{\textbf{99.80}}{99.82} & \pairslash{78.70}{93.76} & \pairslash{93.30}{96.99} & \pairslash{97.40}{99.33} & \pairslash{\underline{91.70}}{96.13} & \pairslash{\underline{91.10}}{97.90} \\
\midrule
\rowcolor[gray]{0.88}
Ours & -- & \pairslash{78.20}{92.05} & \pairslash{\textbf{95.43}}{\textbf{99.19}} & \pairslash{97.77}{\textbf{100.00}} & \pairslash{97.46}{\textbf{100.00}} & \pairslash{\textbf{90.65}}{\textbf{97.70}} & \pairslash{\underline{97.20}}{\textbf{99.92}} & \pairslash{97.53}{\textbf{99.99}} & \pairslash{\textbf{96.45}}{\textbf{99.17}} & \pairslash{\textbf{93.84}}{\underline{98.50}} \\
\bottomrule
\end{tabular}%
}
\end{table*}

\clearpage

\begin{table*}[!t]
\caption{Comparison on UniversalFakeDetect. Each entry reports Acc/AP (\%) for detecting real and fake images from the generator/source named in the row. Mean averages the shared subsets.}
\label{tab:ufd_report}
\centering
\scriptsize
\providecommand{\pair}[2]{#1/#2}
\renewcommand{\arraystretch}{1.03}
\setlength{\tabcolsep}{1.8pt}
\begin{tabularx}{\textwidth}{@{}l*{9}{>{\centering\arraybackslash}X}@{}}
\toprule
Generator & CNNSpot & LGrad & NPR & UnivFD & FreqNet & FatFormer & RINE & C2P-CLIP & Ours \\
Ref. & CVPR~2020 & CVPR~2023 & CVPR~2024 & CVPR~2023 & AAAI~2024 & CVPR~2024 & ECCV~2024 & AAAI~2025 & -- \\
\midrule
ProGAN & \pair{99.99}{100.00} & \pair{99.84}{100.00} & \pair{99.84}{100.00} & \pair{100.00}{100.00} & \pair{97.90}{99.92} & \pair{99.89}{100.00} & \pair{100.00}{100.00} & \pair{99.98}{100.00} & \pair{99.99}{100.00} \\
CycleGAN & \pair{85.20}{93.47} & \pair{85.39}{93.98} & \pair{95.00}{99.53} & \pair{98.50}{98.13} & \pair{95.84}{99.63} & \pair{99.32}{100.00} & \pair{99.30}{100.00} & \pair{97.31}{100.00} & \pair{96.59}{99.83} \\
BigGAN & \pair{70.20}{84.50} & \pair{82.88}{90.69} & \pair{87.55}{94.53} & \pair{94.50}{94.46} & \pair{90.45}{96.05} & \pair{99.50}{99.98} & \pair{99.60}{99.90} & \pair{99.12}{99.96} & \pair{89.33}{99.27} \\
StyleGAN & \pair{85.70}{99.54} & \pair{94.83}{99.86} & \pair{96.23}{99.94} & \pair{82.00}{86.66} & \pair{97.55}{99.89} & \pair{97.15}{99.75} & \pair{88.90}{99.40} & \pair{96.44}{99.50} & \pair{99.34}{100.00} \\
GauGAN & \pair{78.95}{89.49} & \pair{72.45}{79.36} & \pair{86.57}{88.82} & \pair{99.50}{99.25} & \pair{90.24}{99.71} & \pair{99.41}{100.00} & \pair{99.80}{100.00} & \pair{99.17}{100.00} & \pair{82.84}{99.51} \\
StarGAN & \pair{91.70}{98.15} & \pair{99.62}{99.98} & \pair{99.75}{100.00} & \pair{97.00}{99.53} & \pair{93.41}{98.63} & \pair{99.75}{100.00} & \pair{99.50}{100.00} & \pair{99.60}{100.00} & \pair{99.87}{100.00} \\
Deepfake & \pair{53.47}{89.02} & \pair{58.00}{67.91} & \pair{76.89}{84.41} & \pair{66.60}{91.67} & \pair{97.40}{99.92} & \pair{93.23}{97.99} & \pair{80.60}{97.90} & \pair{93.77}{98.59} & \pair{80.35}{94.85} \\
SITD & \pair{66.67}{73.75} & \pair{62.50}{59.42} & \pair{66.94}{97.95} & \pair{63.00}{78.54} & \pair{88.92}{94.42} & \pair{81.11}{97.94} & \pair{90.60}{97.20} & \pair{95.56}{98.92} & \pair{69.44}{99.19} \\
SAN & \pair{48.69}{59.47} & \pair{50.00}{51.42} & \pair{98.63}{99.99} & \pair{57.50}{67.54} & \pair{59.04}{74.59} & \pair{68.04}{81.21} & \pair{68.30}{94.90} & \pair{64.38}{84.56} & \pair{92.24}{97.51} \\
CRN & \pair{86.31}{98.24} & \pair{50.74}{63.52} & \pair{50.00}{50.16} & \pair{59.50}{83.12} & \pair{71.92}{80.10} & \pair{69.45}{99.84} & \pair{89.20}{97.30} & \pair{93.29}{99.86} & \pair{50.02}{99.05} \\
IMLE & \pair{86.26}{98.40} & \pair{50.78}{69.61} & \pair{50.00}{50.16} & \pair{72.00}{91.06} & \pair{67.35}{75.70} & \pair{69.45}{99.93} & \pair{90.60}{99.70} & \pair{93.29}{99.95} & \pair{50.02}{99.90} \\
Guided & \pair{60.07}{73.72} & \pair{77.50}{87.06} & \pair{84.55}{98.26} & \pair{70.03}{79.24} & \pair{86.70}{96.27} & \pair{76.00}{91.99} & \pair{76.10}{96.40} & \pair{69.10}{94.13} & \pair{90.05}{98.29} \\
LDM200 & \pair{54.03}{70.62} & \pair{94.20}{99.03} & \pair{97.65}{99.92} & \pair{94.19}{95.81} & \pair{84.55}{96.06} & \pair{98.60}{99.81} & \pair{98.30}{99.80} & \pair{99.25}{99.99} & \pair{99.30}{99.98} \\
LDM200-CFG & \pair{54.96}{71.00} & \pair{95.85}{99.16} & \pair{98.00}{99.91} & \pair{73.76}{79.77} & \pair{99.58}{100.00} & \pair{94.90}{99.09} & \pair{88.20}{98.30} & \pair{97.25}{99.83} & \pair{97.30}{99.86} \\
LDM100 & \pair{54.14}{70.54} & \pair{94.80}{99.18} & \pair{98.20}{99.92} & \pair{94.36}{95.93} & \pair{65.56}{62.34} & \pair{98.65}{99.87} & \pair{98.60}{99.90} & \pair{99.30}{99.98} & \pair{99.45}{99.99} \\
GLIDE100-27 & \pair{60.78}{80.65} & \pair{87.40}{93.23} & \pair{96.25}{99.87} & \pair{79.07}{93.93} & \pair{85.69}{99.80} & \pair{94.35}{99.13} & \pair{88.90}{98.80} & \pair{95.25}{99.72} & \pair{99.45}{99.96} \\
GLIDE100-10 & \pair{63.80}{84.91} & \pair{90.70}{95.10} & \pair{97.15}{99.89} & \pair{79.85}{95.12} & \pair{97.40}{99.78} & \pair{94.65}{99.41} & \pair{92.60}{99.30} & \pair{95.25}{99.79} & \pair{99.15}{99.90} \\
GLIDE50-27 & \pair{65.66}{82.07} & \pair{89.55}{94.93} & \pair{97.35}{99.92} & \pair{78.14}{94.59} & \pair{88.15}{96.39} & \pair{94.20}{99.20} & \pair{90.70}{98.90} & \pair{96.10}{99.83} & \pair{99.35}{99.95} \\
DALL-E & \pair{55.58}{70.59} & \pair{88.35}{97.23} & \pair{87.15}{99.26} & \pair{86.78}{88.45} & \pair{59.06}{77.78} & \pair{98.75}{99.82} & \pair{95.00}{99.30} & \pair{98.55}{99.91} & \pair{93.95}{99.61} \\
\midrule
\rowcolor[gray]{0.88}
Mean & \pair{69.58}{83.58} & \pair{80.28}{86.35} & \pair{87.56}{92.76} & \pair{81.38}{90.14} & \pair{85.09}{91.95} & \pair{90.86}{98.16} & \pair{91.31}{98.78} & \pair{93.79}{98.66} & \pair{88.84}{99.30} \\
\bottomrule
\end{tabularx}
\end{table*}

Tables~\ref{tab:genimage_report} and~\ref{tab:ufd_report} expand the external benchmark results to generator-level Acc/AP comparisons. On GenImage, Ours gives the highest mean Acc among the retained baselines with complete per-generator Acc/AP, while BigGAN remains the main fixed-threshold weakness under the SD v1.4-source calibration. On UniversalFakeDetect, Ours keeps the highest mean AP among the listed methods on the shared 19-subset comparison, but its fixed-threshold mean Acc is lower than C2P-CLIP, RINE, and FatFormer. This again separates ranking quality from the deployed threshold and is consistent with the 21-generator summary in Table~\ref{tab:benchmark_summary}.

\clearpage

\subsection{Component Ablation}
\label{sec:ablation_study}

\begin{table}[t]
\caption{Core Component Ablation on AIGCDetectBenchmark}
\label{tab:prior_ablation}
\centering
\scriptsize
\renewcommand{\arraystretch}{1.08}
\setlength{\tabcolsep}{4pt}
\begin{tabular}{lrrrr}
\toprule
Variant & Acc & $\Delta$ & Real & Fake \\
\midrule
Full FreqPRISM & 94.82 & 0.00 & 94.69 & 94.95 \\
No artifact & 84.06 & -10.76 & 96.62 & 71.50 \\
No semantic & 91.62 & -3.20 & 89.17 & 94.08 \\
No residual & 90.62 & -4.21 & 93.23 & 88.00 \\
No native tile & 93.05 & -1.78 & 95.55 & 90.54 \\
\midrule
Artifact only & 85.62 & -9.20 & 78.56 & 92.69 \\
Semantic only & 78.73 & -16.10 & 94.70 & 62.75 \\
Residual only & 88.99 & -5.83 & 90.99 & 86.99 \\
\bottomrule
\end{tabular}
\end{table}

Table~\ref{tab:prior_ablation} identifies which priors drive the fixed-threshold result. Removing artifact evidence causes the largest loss, reducing fake accuracy from 94.95 to 71.50. This makes the artifact prior the main fake-side anchor. Removing semantic evidence mainly harms real-side calibration, lowering Real accuracy from 94.69 to 89.17. Removing residual evidence reduces Acc by 4.21 points and fake accuracy by 6.95 points, showing that the residual branch contributes non-redundant correction even though it is used with a low logit weight. Removing native tiles has the smallest mean-Acc effect, but it lowers fake accuracy by 4.41 points; the tile path therefore mainly improves recall for local high-resolution evidence.

Additional diagnostics are consistent with this interpretation. In the artifact-family ablation, removing codec/block evidence lowers mean Acc to 84.97 and fake accuracy to 70.25, whereas the codec/block-only variant remains close to the full artifact setting. In the tile diagnostic, the current native top-1 tile policy outperforms whole-image-only scoring by 1.78 Acc points. In the residual sweep, the selected residual scale improves over no residual while avoiding the larger ranking loss observed at a 2.00 gamma scale. These results support the design choice of combining the priors as constrained evidence sources rather than treating any single branch as a complete detector.

\section{Discussion}
\label{sec:discussion}
Discuss limitations, broader implications, and possible future improvements.

\section{Conclusion}
\label{sec:conclusion}
Summarize the main findings and contributions.

\bibliographystyle{IEEEtran_doi}
\bibliography{ref,uncited}

\clearpage
\section*{Reproducibility Checklist}

Table~\ref{tab:aigcdetect_ap_comparison} reports the available per-generator AP values for AIGCDetectBenchmark. It follows the same generator order and highlighting convention as Table~\ref{tab:aigcdetect_comparison}.

\setcounter{table}{4}
\begin{table*}[!t]
\caption{Comparison on AIGCDetectBenchmark. Each entry reports AP (\%) for detecting real and fake images from the generator named in the column. DIRE-G denotes the ProGAN-source DIRE variant, whereas DIRE-D follows the ADM-source DIRE setup; SDAIE is the one-class variant, whereas SDAIE$^\dagger$ is the binary ProGAN-source variant. The best and second-best values in each column are marked in bold and underline, respectively.}
\label{tab:aigcdetect_ap_comparison}
\centering
\scriptsize
\providecommand{\rotgen}[1]{\makebox[0pt][l]{\rotatebox{60}{\textit{#1}}}}
\renewcommand{\arraystretch}{1.03}
\setlength{\tabcolsep}{2.0pt}
\begin{tabularx}{\textwidth}{@{}l>{\centering\arraybackslash}p{0.072\textwidth}*{18}{Y}@{}}
\toprule
Method & Ref. & \rotgen{ProGAN} & \rotgen{StyleGAN} & \rotgen{BigGAN} & \rotgen{CycleGAN} & \rotgen{StarGAN} & \rotgen{GauGAN} & \rotgen{StyleGAN2} & \rotgen{WFIR} & \rotgen{ADM} & \rotgen{GLIDE} & \rotgen{Midjourney} & \rotgen{SD v1.4} & \rotgen{SD v1.5} & \rotgen{VQDM} & \rotgen{Wukong} & \rotgen{DALL-E 2} & \rotgen{SDXL} & \rotgen{Mean} \\
\midrule
CNNSpot & CVPR~2020 & \textbf{100.0} & 99.8 & 86.0 & 94.9 & 99.0 & 90.8 & 99.5 & \textbf{99.9} & 75.7 & 72.3 & 66.2 & 61.2 & 61.6 & 68.8 & 57.3 & 53.5 & 72.6 & 79.95 \\
FreDect & ICML~2020 & \textbf{100.0} & 89.0 & 93.6 & 84.8 & 99.5 & 82.8 & 82.5 & 55.9 & 61.8 & 52.9 & 46.1 & 37.8 & 37.8 & 85.1 & 39.6 & 38.2 & 49.5 & 66.87 \\
Fusing & ICIP~2022 & \textbf{100.0} & 99.5 & 90.7 & 95.5 & \underline{99.8} & 88.3 & 99.6 & 93.3 & 94.1 & 77.5 & 70.0 & 65.4 & 65.7 & 75.6 & 64.6 & 68.1 & 79.4 & 83.95 \\
GramNet & CVPR~2020 & \textbf{100.0} & 99.2 & 81.8 & 95.3 & 99.2 & 85.0 & 99.1 & 95.2 & 73.1 & 66.8 & 56.8 & 59.8 & 60.4 & 61.1 & 55.6 & 49.8 & 68.2 & 76.86 \\
LNP & ECCV~2022 & \textbf{100.0} & 99.3 & 94.5 & 89.5 & \textbf{100.0} & 84.5 & 99.7 & 42.8 & 93.4 & 92.8 & 86.9 & 96.3 & 96.0 & 94.9 & 95.3 & 98.3 & 87.8 & 91.29 \\
LGrad & CVPR~2023 & \textbf{100.0} & 98.3 & 92.9 & 95.0 & \textbf{100.0} & 95.4 & 97.9 & 58.0 & 73.0 & 80.4 & 71.9 & 62.4 & 62.9 & 77.5 & 62.5 & 82.6 & 80.0 & 81.80 \\
DIRE-G & ICCV~2023 & \underline{99.1} & 91.7 & 75.3 & 80.6 & 99.3 & 72.2 & 88.3 & 60.1 & 85.8 & 78.4 & 61.9 & 49.9 & 49.5 & 54.6 & 55.4 & 74.5 & 54.0 & 72.38 \\
DIRE-D & ICCV~2023 & 58.8 & 56.7 & 46.9 & 50.0 & 40.6 & 47.3 & 58.0 & 59.0 & \textbf{99.8} & 99.5 & 97.3 & 98.6 & 98.8 & 99.0 & 98.4 & 99.7 & 99.1 & 76.92 \\
UnivFD & CVPR~2023 & \textbf{100.0} & 97.6 & 99.3 & 99.8 & 99.4 & \textbf{100.0} & 97.9 & 96.7 & 86.8 & 83.8 & 74.0 & 86.1 & 85.8 & 96.5 & 91.1 & 63.0 & 67.6 & 89.73 \\
PatchCraft & arXiv~2023 & \textbf{100.0} & 99.0 & 99.4 & 85.3 & \textbf{100.0} & 81.3 & 97.7 & 95.3 & 93.4 & 94.0 & 96.5 & 99.1 & 99.1 & 96.3 & 97.5 & 99.6 & \underline{99.9} & 96.07 \\
NPR & CVPR~2024 & \textbf{100.0} & \underline{99.9} & 94.0 & 99.1 & \textbf{100.0} & 88.7 & \textbf{100.0} & 94.0 & 98.3 & \underline{99.6} & 99.0 & 99.7 & 99.6 & 99.2 & 98.9 & \textbf{100.0} & 99.6 & 98.21 \\
AIDE & ICLR~2025 & \textbf{100.0} & \textbf{100.0} & 94.4 & \underline{99.9} & \textbf{100.0} & 97.7 & \textbf{100.0} & \underline{99.6} & 99.0 & 99.1 & 88.4 & 98.2 & 98.2 & \underline{99.4} & 98.6 & 99.5 & 99.4 & 98.31 \\
SDAIE & TPAMI~2026 & 87.3 & 89.7 & 66.0 & 85.9 & 99.1 & 73.4 & 89.5 & 79.4 & 95.9 & 99.3 & 92.8 & 96.5 & 96.4 & 97.2 & 95.6 & 95.9 & 96.6 & 90.39 \\
SDAIE$^\dagger$ & TPAMI~2026 & \textbf{100.0} & \textbf{100.0} & 97.3 & 99.4 & \textbf{100.0} & 89.9 & \textbf{100.0} & 98.1 & 98.3 & 98.5 & 99.4 & \underline{99.9} & \underline{99.8} & 97.9 & 99.4 & \underline{99.9} & \textbf{100.0} & \underline{98.69} \\
FatFormer & CVPR~2024 & \textbf{100.0} & 99.8 & \textbf{100.0} & \textbf{100.0} & \textbf{100.0} & \textbf{100.0} & \underline{99.9} & 98.5 & 91.7 & 96.0 & 62.8 & 81.1 & 81.1 & 97.0 & 85.9 & 81.8 & 86.9 & 91.90 \\
RINE & ECCV~2024 & \textbf{100.0} & 99.4 & \underline{99.9} & \textbf{100.0} & \textbf{100.0} & \textbf{100.0} & \textbf{100.0} & 99.5 & 96.3 & 97.9 & 87.7 & 98.4 & 98.3 & 99.2 & 98.6 & 89.6 & 96.3 & 97.71 \\
C2P-CLIP & AAAI~2025 & 83.5 & 65.8 & 79.0 & 98.3 & 87.8 & 71.3 & 66.5 & 99.5 & 97.4 & 98.8 & \underline{99.5} & \textbf{100.0} & \textbf{99.9} & 98.8 & \textbf{99.7} & 99.8 & \underline{99.9} & 90.91 \\
D$^3$ & CVPR~2025 & \textbf{100.0} & 99.2 & \textbf{100.0} & 98.5 & 99.4 & \underline{99.9} & 99.4 & 91.6 & \underline{99.5} & 99.5 & \underline{99.5} & 99.8 & 99.7 & \textbf{99.6} & \underline{99.6} & 86.5 & 98.8 & 98.25 \\
CO-SPY & CVPR~2025 & \textbf{100.0} & \textbf{100.0} & 98.2 & 99.4 & \textbf{100.0} & 98.1 & \textbf{100.0} & 80.2 & 89.6 & 97.2 & 94.3 & 93.2 & 93.0 & 91.2 & 89.6 & 97.2 & 98.2 & 95.25 \\
\midrule
\rowcolor[gray]{0.88}
Ours & -- & \textbf{100.0} & \textbf{100.0} & 99.3 & 99.8 & \textbf{100.0} & 99.5 & \textbf{100.0} & 97.5 & \underline{99.5} & \textbf{99.9} & \textbf{99.6} & 98.6 & 98.8 & 99.1 & 97.5 & \underline{99.9} & \underline{99.9} & \textbf{99.34} \\
\bottomrule
\end{tabularx}
\end{table*}

\end{document}
